% 01-description.tex - Description and Scope

\chapter{Description}
\label{sec:description}

\section{Purpose}

This document defines the interface between the Data Recorder (DR) subsystem and
other Long Wavelength Array (LWA) station subsystems.  The DR records the output
of the Next-generation Digital Processor (NDP) subsystem and is controlled by the
Monitor and Control System (MCS).  Whereas station architecture and subsystem ICDs
may refer to DR as a whole, this ICD applies to a single DR instance.

\begin{newInVersion}{2.4.1}
In previous versions of this document, the upstream digital processor was referred
to as DP (Digital Processing).  All references have been updated to NDP
(Next-generation Digital Processor) to reflect the current system architecture.
\end{newInVersion}

\section{Scope}

This ICD describes the DR's physical and electrical connections, software
interfacing, and control methods.  It covers:
\begin{itemize}
    \item Physical and network connections between DR and other station subsystems.
    \item The MCS message-based command and monitoring interface.
    \item Configuration parameters and system requirements.
    \item Data format specifications for recorded and spectrometer output files.
\end{itemize}

\section{System Overview}

The DR runs the Data Recorder Operating Software version~2 (DROS2), which
provides all ICD-specified functionality including MCS message processing,
data recording, spectrometer computation, and storage management.

The build system compiles two DROS2 binaries from the same source code:

\begin{enumerate}
    \item \textbf{Spectrometer mode} (primary): The DR records raw NDP data
          streams and can compute integrated power spectra of beam data in real
          time using the \cmd{SPC} command.  This is the standard deployment mode.
    \item \textbf{LiveBuffer mode} (secondary): The DR provides a live capture
          buffer for on-the-fly packet capture via the \cmd{BUF} command and
          \mib{BUFFER} MIB entries.  This mode is not the standard deployment.
\end{enumerate}

The active mode is a runtime selection controlled by the contents of the
\texttt{/LWA/config/binary-mode} file, which specifies either
\texttt{Spectrometer} or \texttt{LiveBuffer}.  The corresponding binary is
launched when the DROS2 service starts; switching modes requires only updating
this file and restarting the service.
DROS2 is deployed as a \texttt{systemd} service (\texttt{dros.service}) that
starts automatically on boot, restarts on failure, and logs to the system
journal.

\section{Multi-DR Deployment}
\label{sec:multi-dr}

\begin{newInVersion}{2.4.1}
A station may deploy multiple DR instances on a single host, each recording from
one or more NDP data streams.  Multi-DR deployment is the preferred configuration
and is managed through the \texttt{build\_helper.py} utility, which creates
per-instance build directories with separate:

\begin{itemize}
    \item Installation paths (\texttt{/LWA/DR\textit{N}/})
    \item Storage mount points (\texttt{/LWA\_STORAGE/DR\textit{N}/})
    \item Configuration files (\texttt{/LWA/DR\textit{N}/config/})
    \item Runtime directories (\texttt{/LWA/DR\textit{N}/runtime/})
    \item Service files (\texttt{dros-dr\textit{N}.service})
\end{itemize}

Each DR instance has its own \texttt{defaults\_v2.cfg} configuration file
specifying a unique reference designator (e.g., \texttt{DR1}, \texttt{DR2}),
IP address, and port numbers.  When running in a single-DR configuration, the
default installation paths (\texttt{/LWA/}, \texttt{/LWA\_STORAGE/}) are used
directly.
\end{newInVersion}

\section{Document Organization}

The remainder of this document is organized as follows:

\begin{description}
    \item[Section~\ref{sec:abbreviations}] defines abbreviations, acronyms, and
        conventions used throughout this document.
    \item[Section~\ref{sec:physical}] describes the physical and electrical
        interfaces of the DR.
    \item[Section~\ref{sec:monitor-control}] describes the monitor and control
        interface, including configuration, MIB structure, and control commands.
    \item[Section~\ref{sec:safety}] describes the safety monitoring interface.
    \item[Section~\ref{sec:references}] lists references and related documents.
    \item[Appendix~\ref{app:spectrometer}] defines the spectrometer data file
        format.
    \item[Appendix~\ref{app:errors}] lists error messages.
    \item[Appendix~\ref{app:data-formats}] provides a reference of all supported
        data formats.
\end{description}
